
\section{Results}

\subsection{Method Comparison}

The results of the Leiden and Louvain algorithms applied to the Bac and Vir networks are summarized in the following table:

\begin{table}[htbp]
\centering
\caption{Community detection results for bacterial and viral networks.}
\label{tab:community_results}
\begin{tabular}{lcccc}
\hline
Dataset & Method & \# Communities & Modularity & Runtime (s) \\
\hline
Bac & Louvain & 5  & 0.374 & 0.340 \\
Bac & Leiden  & 6  & 0.333 & 0.017 \\
Vir & Louvain & 23 & 0.892 & 0.136 \\
Vir & Leiden  & 23 & 0.884 & 0.007 \\
\hline
\end{tabular}
\end{table}

For the Bac network, the Louvain algorithm identified 5 communities, 
whereas the Leiden algorithm detected 6 communities.
The slight increase in the number of communities produced by Leiden suggests a 
refinement of the partition obtained by Louvain. This behavior is consistent with the theoretical properties
of the Leiden algorithm, which aims to ensure better-connected and more internally coherent communities.

In terms of modularity, Louvain achieved a value of approximately 0.374,
while Leiden reached approximately 0.333. 
The higher modularity value for Louvain in the bacterial network suggests that 
Louvain is more effective at optimizing the modularity objective in this specific case.
This indicates that, although Leiden produces a slightly more refined partition, 
the Louvain partition achieves stronger optimization of the modularity objective under the given settings.

For the Vir network, which has approximately twice as many nodes as the Bac network,
both Louvain and Leiden identified 23 communities in the viral network.
The strong agreement between the two methods indicates that the underlying community 
structure is clearly defined and does not depend strongly on the specific optimization strategy.

In terms of modularity, Louvain achieved a value of approximately 0.892,
while Leiden obtained approximately 0.884.
Although Louvain reaches a slightly higher modularity score,
the difference between the two methods is small. 
This indicates that both algorithms produce partitions of very similar quality 
with respect to modularity optimization.

As expected, the runtime of the Louvain algorithm was higher than that of the Leiden algorithm in both networks.
For the Bac network, Louvain required approximately 0.340 seconds, while Leiden required approximately 0.017 seconds.
For the Vir network, Louvain required approximately 0.136 seconds, while Leiden required approximately 0.007 seconds.

This difference in runtime can be attributed to more efficient local moving, a better-structured
aggregation phase, and faster convergence of the Leiden algorithm compared to Louvain, 
especially in larger networks.

\subsection{ARI}

The Adjusted Rand Index (ARI) further confirms the strong agreement between Louvain and Leiden. 
For the viral network, the ARI is approximately 0.979, 
indicating an almost identical partitioning of the network by both algorithms.

For the bacterial network, the ARI is slightly lower at approximately 0.958,
yet it still reflects a very high degree of similarity. 
The minor difference can be attributed to the additional community detected by Leiden, 
which suggests a small refinement of the Louvain partition rather than a 
fundamentally different community structure.

Overall, the high ARI values demonstrate that both methods 
produce highly consistent community assignments across both datasets.

\subsection{Stability of Communities}

The observed differences between the Bac and Vir networks can be attributed to their fundamentally
different structural organization. The bacterial network is denser and more strongly connected,
indicating the presence of tightly interconnected protein communities.
Such dense connectivity implies that many proteins are linked across different regions of
the network, resulting in a cohesive structure with numerous inter-community connections.
In contrast, the viral network is sparser and contains a larger proportion of weakly
connected nodes. Connections between modules are less frequent, leading to sharper boundaries
between communities.
These structural differences are directly reflected in the modularity values. Modularity
measures the extent to which connections are concentrated within communities rather than 
between them.
In the dense bacterial network, the abundance of cross-community edges reduces
the contrast between intra- and inter-community connectivity. As a result, the modularity 
remains comparatively lower.
In the sparse viral network, the clear separation between communities leads to a higher modularity score,
as the majority of connections are concentrated within communities rather than between them.
Thus, the difference in modularity between the two datasets does not indicate better algorithmic 
performance on the viral network, but rather reflects the intrinsic topological differences between 
a densely interconnected graph and a more clearly segmented modular structure.




